\begingroup%
\definecolor{ksBlue}{HTML}{3478B8}%
\definecolor{ksTeal}{HTML}{27877F}%
\definecolor{ksAmber}{HTML}{C58B28}%
\definecolor{ksInk}{HTML}{243447}%
\definecolor{ksPaleBlue}{HTML}{F0F6FC}%
\definecolor{ksPaleTeal}{HTML}{EFF8F5}%
\begin{tikzpicture}[
  x=1mm, y=1mm,
  card/.style={draw, fill=white, rounded corners=2mm, line width=.7pt,
    text=ksInk, align=flush center, font=\fontsize{9.5}{12.5}\selectfont,
    inner xsep=2mm, inner ysep=1.7mm, outer sep=.35pt},
  future/.style={card, draw=ksBlue, minimum width=69mm},
  current/.style={card, draw=ksTeal},
  frame/.style={draw, rounded corners=3mm, line width=.7pt, inner sep=3mm,
    outer sep=.35pt},
  heading/.style={text=ksInk, font=\bfseries\fontsize{10}{12.5}\selectfont,
    inner sep=0pt, outer sep=0pt},
  flow/.style={-{Stealth[length=2.25mm,width=1.65mm]}, line width=.7pt,
    rounded corners=2mm}
]
  \node[heading] (future-title) at (83,0) {面向未来：可验证的自动化闭环};
  \node[heading] (now-title) at (0,0) {面向当下：MaaS 推理的算子需求};

  \node[future] (f1) at (83,-14) {模型能力提升的新来源\\
    RL 合成数据 · 系统自我迭代 (RSI)};
  \node[future] (f2) at (83,-32) {前提：边界清楚、\\
    结果能客观验证的好问题};
  \node[future] (f3) at (83,-53) {算子开发天然满足这一前提\\
    正确性可判定（编译 + 数值对齐）\\
    性能反馈连续（延迟 / Roofline）\\
    软硬件协同的组合优化，有足够难度};

  \node[current, minimum width=38.5mm, minimum height=23mm] (n1) at (-20.5,-20)
    {模型架构快速演进\\
    Dense / MoE / MLA\\
    线性与稀疏注意力\\
    FP8 / MXFP8 / NVFP4};
  \node[current, minimum width=38.5mm, minimum height=23mm] (n2) at (20.5,-20)
    {异构推理卡多代并存\\
    Hopper / Blackwell\\
    Ascend 等国产卡\\
    HIP 等其他平台};
  \node[current, minimum width=75mm] (n3) at (0,-43)
    {算子需求呈 M × N × K 组合爆炸\\
    HPC 专家少，手工调优周期长};
  \node[current, minimum width=75mm] (n4) at (0,-59)
    {算子效率直接影响\\
    单 Token 成本 · 算力利用率 · 服务延迟};

  \coordinate (future-base) at (83,-65);
  \coordinate (now-base) at (0,-65);
  \begin{scope}[on background layer]
    \node[frame, draw=ksBlue, fill=ksPaleBlue,
      fit=(future-title)(f1)(f2)(f3)(future-base)] (future-group) {};
    \node[frame, draw=ksTeal, fill=ksPaleTeal,
      fit=(now-title)(n1)(n2)(n3)(n4)(now-base)] (now-group) {};
  \end{scope}

  \node[card, draw=ksAmber, line width=.85pt, minimum width=96mm] (core)
    at (41.5,-87) {目标：构建可验证、能收敛、可理解的\\
    自动化算子生成闭环，并把方法用到更多任务上};
  \coordinate (future-turn) at (62,-76);
  \coordinate (now-turn) at (21,-76);

  \draw[flow, ksBlue] (f1) -- (f2);
  \draw[flow, ksBlue] (f2) -- (f3);
  \draw[flow, ksTeal] (n1.south) -- (n1.south |- n3.north);
  \draw[flow, ksTeal] (n2.south) -- (n2.south |- n3.north);
  \draw[flow, ksTeal] (n3) -- (n4);
  \draw[flow, ksBlue] (future-group.south) |- (future-turn) --
    (future-turn |- core.north);
  \draw[flow, ksTeal] (now-group.south) |- (now-turn) --
    (now-turn |- core.north);
\end{tikzpicture}%
\endgroup%
